\section{Composing signed horizons with oriented power pools}

The reused CRT anchor and the oriented power theorem have complementary
strengths. The former removes a derived-label range request; the latter
combines the remaining power certificates. Their composition requires a new
sign argument, because the factorial normalization used in the earlier
pooling proof is absent after the anchor reduction.

\subsection{The common arithmetic interface}

The five bases are
\[
 n_0=t,\qquad n_1=B,\qquad n_2=BH,\qquad n_3=q,\qquad n_4=q+1,
\]
with respectively $1,1,1,3,10$ power requests. The retained canonical signed
Pell data at center $A$ satisfy
\[
 |A|>W^W,\qquad n_j^2,m_i^2,k_i^2<W,\qquad W\ge258.
\]
This interface is established before the signs of the arithmetic coordinates
are decoded. Every requested index has a canonical signed Pell pair
$(X_i,Y_i)$, with $D_i=X_i-AY_i$.

Define, for each block,
\begin{align*}
 T_j&=2An_j-n_j^2-1,\\
 E_j&=\sum_{i\in j}\bigl((m_i-D_i)^2-Y_i^2\bigr)
       -2n_j\sum_{i\in j}Y_i(m_i-X_i).
\end{align*}
The five-base residual is the ordinary polynomial
\[
 R=\sum_{j=0}^4 n_jE_j\prod_{\ell\ne j}T_\ell
       -z\prod_{j=0}^4T_j.
\]
The inherited oriented-pooling theorem makes $R=0$ equivalent to exactness
of every request at a nonzero actual base. Zero-base blocks are temporarily
unconstrained, and are excluded by the argument below. For nonzero bases,
$\operatorname{sgn}(n_j/T_j)=\operatorname{sgn}(A)$; this common sign is what
prevents cancellation of the squared power errors.

Two smaller-degree alternatives use three and two quotient coordinates.
The paired version retains the direct $t$ singleton, pools $B,BH$ with base
weights, and pools $q,q+1$ without weights. The staged version pools
$t,B,BH$ with base weights and again uses the unweighted $q,q+1$ pool.
All final equations remain sums of squares of polynomial residuals.

\subsection{A parity argument replaces the removed factorial certificate}

For the moderate-coefficient source set
\[
 t=128+256u,\qquad K=t^4,
\]
and retain the request $F=t^K$. The smaller-source alternative sets
\[
 t=16+32u,\qquad K=2^{20}t.
\]
All parameters here range over all integers. In both cases $t$ is a nonzero
multiple of four. In the first case $K$ is positive and even; in the second,
the signed-power theorem forces $K\ge0$ and thus $t>0$. In either case the
first power gives a positive even $F$ with $F>|t|$.

The retained CRT row reads
\[
 Bv=C+F+1,\qquad C=2c.
\]
Its right-hand side is odd. Consequently $B$ is odd and nonzero, independently
of its eventual binomial interpretation. Also $H=2+4S$ for an integer $S$,
so $|BH|\ge2$. Its exact power gives
\[
 x\ge0,\qquad q=(BH)^x.
\]
Hence $q\ne0,-1$, and all five bases are nonzero. This proves that every
one of the old power requests is recovered by the new pools. It also justifies
the unweighted two-base variants, because the nonzero consecutive integers
$q$ and $q+1$ have the same sign.

\subsection{The endpoint removes the remaining negative horizon}

For the moderate-coefficient source, the sign of $t$ has not yet been
recovered. The retained beta endpoint rows give
\[
 y=F(2v+1),\qquad g=1+(t+1)y,\qquad x=1+g u_t.
\]
Since $|t|\ge128$ and $|y|\ge F$,
\[
 |g|\ge |t+1|F-1\ge127F-1>1.
\]
Thus $x=0$ would force $g u_t=-1$, a contradiction. It follows that $x\ge1$
and $|q|\ge2$. The literal retained request is $q^t$, with nominal index
$t$; its signed exactness forces $t\ge0$. Therefore $t\ge128$.

This argument uses neither a missing factorial row nor a false conclusion
from the binomial high quotient. In particular, setting $q=1$ in the digit
row does not identify $B$ alone: the high quotient remains. The endpoint
factorization is the required noncircular exclusion of $x=0$.

\subsection{Completeness and literal polynomial projections}

The positive branches of both affine progressions are cofinal and divisible
by four. Package 173's completeness construction can choose its horizon in
either progression. It needs precisely these properties and
$K\ge2^{20}t$, rather than the earlier square representation. The CRT anchor
$C$ can still be chosen positive, even, and larger than every other nominal
index. All subsequent cap and arithmetic witnesses can be completed.

Conversely the preceding sign deductions restore the hypotheses of the
retained arithmetic decoding theorem. The same label-separation and Horn-flow
arguments yield a finite C222 refutation, and the identified inherited
C222/ZF/ZFC calibration gives a finite ZFC refutation. Thus the new sources
have a zero exactly when the fixed ZFC presentation has a finite refutation.

Every subsequent projection uses a residual
$\epsilon v+f=0$ with $\epsilon\in\{1,-1\}$ and $v$ absent from $f$.
The supplied definition lists are acyclic. Substitution therefore gives a
bijection on integer zero sets, with the integral polynomial reconstruction
$v=-\epsilon f$. Registers in the final circuits are abbreviations and do
not add quantified variables.

The retained profiles are as follows. The last column bounds literal circuit
constants, not fully expanded polynomial coefficients.
\begin{center}
\begin{tabular}{rrrrr}
\hline
Variables & Exact degree & Residuals & Operations & Literal constant bound\\
\hline
131&12&70&2475&1099511627780\\
132&12&71&2484&256\\
117&16&56&2439&256\\
108&32&47&2434&256\\
97&128&36&2418&256\\
91&256&30&2400&256\\
87&32768&26&2388&256\\
86&2097780&25&2385&256\\
\hline
\end{tabular}
\end{center}

The immediate Oriented Powers comparisons include $137$ variables at degree
$12$, $121$ at degree $16$, $111$ at degree $32$, $99$ at degree $128$, and
$87$ at degree $8388612$. These are matched construction improvements, not
claims of global minimality. Exact degrees are certified by nonzero leading
homogeneous specializations; propagated coefficient-norm estimates provide
separate explicit upper bounds on expanded coefficient height.

The source checker independently binds all sixteen actual requests, fourteen
canonical Pell pairs, their shifted multiplication wires and norms, the
complete orientation row, and every squared frame atom. The circuit checker
verifies local polynomial expansions, acyclic unit pivots, source-row
coverage, degree certificates, and a nonzero derivative witness for each
quantified variable. These finite checks authenticate the construction;
they do not replace its unbounded arithmetic proof or the inherited logical
calibration.
